\documentclass[a4paper,openany]{book} % DIN A4 Buch, Kapitelbeginn auch auf geraden Seiten (links)
\usepackage[T1]{fontenc}
\usepackage[utf8]{inputenc}

% f"ur Trennung der W"orter nach neuer dt. Rechtschreibung
\usepackage{ngerman}

% Grafik packages
\usepackage{lscape} % f"ur die Verwendung der Landscape-Umgebung
\usepackage{color} % f"ur farbigen Text
\usepackage[dvips]{graphicx}
\usepackage{picinpar} % f"ur Text um Grafik
\usepackage{picins} % f"ur Text um Grafik
\usepackage{wrapfig}

% Seitengr"o"se
\renewcommand{\evensidemargin}{0in} % Abstand nach au"sen bei ungeraden Seiten
\renewcommand{\oddsidemargin}{0in} % Abstand nach au"sen bei geraden Seiten
\renewcommand{\textwidth}{16cm} % Breite des Textblocks ohne Kopf- und Fu"szeile
\renewcommand{\textheight}{22cm} % Hhe des Textblocks ohne Kopf- und Fu"szeile
\renewcommand{\footskip}{2cm}

% f"ur die Verwendung von mehrspaltigen Texten
\usepackage{multicol}

% f"ur die Indexerstellung/-darstellung
\usepackage{makeidx}
\usepackage{index}
  \newindex{default}{idx}{ind}{Schlagwortverzeichnis}
  \newindex{personen}{pdx}{pnd}{Personenverzeichnis}
  \newindex{regionen}{rdx}{rnd}{Städte- und Regionsverzeichnis}

% \usepackage{showidx} % einbinden, um indizierte Worte auf den Seiten anzuzeigen
% Pakete zur Verwaltung der Kopf- und Fu"szeilen
%\usepackage{fancyheadings}
\usepackage{fancyhdr}
\pagestyle{fancyplain}

% Einfaches Einbinden von EPS-Dateien
\usepackage{epsfig}

% Zugriff auf Anzahl Seiten und
% Seitennummer der letzten Seite
\usepackage{lastpage} % \pageref{LastPage} -> Seitennummer der letzten Seite
\usepackage{totpages} % \ref{TotPages} -> Gesamtzahl Seiten

% Kopfzeilendefinition
% [] f"ur gerade Seiten
% {} f"ur ungerade Seiten
\lhead[\let\uppercase\relax\leftmark]{\let\uppercase\relax\rightmark} % links-oben
\chead{} % mitte-oben
\rhead[\let\uppercase\relax\rightmark]{\let\uppercase\relax\leftmark} % rechts-oben

% Fu"szeilendefinition
% [] f"ur gerade Seiten
% {} f"ur ungerade Seiten
\lfoot[Seite \thepage]{\today} % links-unten
\cfoot{Drachend{\ae}mmerung - eine Rollenspielwelt} % mitte-unten
\rfoot[\today]{Seite \thepage} % rechts-unten

% Definition der Kopfzeilenlinien
\renewcommand{\headrulewidth}{1pt}
\renewcommand{\plainheadrulewidth}{1pt}

% Definition der Fuï¿œeilenlinien
\renewcommand{\footrulewidth}{1pt}
\renewcommand{\plainfootrulewidth}{1pt}

% Groï¿œzï¿œgiges Verhalten beim Zeilenumbruch
% Lï¿œsst auch groï¿œe Lï¿œcken zwischen Worten zu
%\sloppy

% Erweiterte Tabellenumgebungen
\usepackage{array}
\usepackage{tabularx}
\usepackage{longtable}